\section{Background}
\label{sec:background}

\subsection{Medical foundation models}
\label{sec:models}

%% [written after the `models' literature bucket is in; see notes/lit_models.md]
\input{sections/sec2_1_models}

\subsection{Interpretability preliminaries and notation}
\label{sec:prelim}

\textbf{Terminology.} \emph{Feature space}, \emph{representation} and \emph{internal representation}
are synonyms here, and \emph{feature-space interpretability} means measuring or intervening on them.
Beyond that the words \emph{interpretability} and \emph{explainability} are not distinguished, as in
several surveys~\cite{wan2025survey}; the properties they are sometimes used to separate, faithfulness, plausibility and intrinsic transparency, are distinguished explicitly where they matter
(\S\ref{sec:evaluation})~\cite{lipton2016mythos,rudin2018stop,rudin2021interpretable}.
\emph{Post-hoc} methods analyse an unchanged trained model; \emph{intrinsic} designs make part of its
computation interpretable by construction~\cite{rudin2018stop,rudin2021interpretable}.
A \emph{local} explanation concerns one input, a \emph{global} explanation a model or population;
feature-space interpretability is mostly global and post hoc. Input-attribution methods (saliency, class activation maps~\cite{selvaraju2017gradcam}, LIME~\cite{ribeiro2016why},
SHAP~\cite{lundberg2017unified}, integrated gradients~\cite{sundararajan2017axiomatic}, layer-wise relevance
propagation~\cite{bach2015pixel}) return per-input maps silent about the representation and enter only
where they test a representation-level claim (\S\ref{sec:evaluation}).

\textbf{Notation.} A medical foundation model receives an image, volume or slide $x$ and optionally a
text $t$. A vision encoder $f_v$ produces patch-level features and a pooled embedding; a dual encoder
adds a text encoder $f_t$, both mapped into a shared space. In a multimodal LLM, a connector $\pi$
produces visual tokens and the language model hidden states $h^{(\ell)}$ at every layer and position. We
write $z_\lambda(x)$ for the representation at locus $\lambda$ (\S\ref{sec:loci}), $c$ for a concept
with externally defined value $c(x)$ and $g$ for a decoder from $z_\lambda$ to $c$. An intervention
replaces $z_\lambda$ with $\tilde z_\lambda$ and measures a pre-specified behavioural functional $B$ (a
logit, probability, report-quality score or downstream metric), with generated text mapped to $B$ and
decoding fixed or averaged.

\textbf{Two questions, two kinds of evidence.} \emph{What does the representation encode?} is addressed
observationally, by fitting $g$, computing similarities, decomposing $z_\lambda$ and locating information
along the path; \emph{what does the model use?} by intervening on $z_\lambda$ and observing $B$. The
two are not interchangeable: a concept can be linearly decodable from a representation the output never
uses, so probing accuracy overstates use, as control tasks~\cite{hewitt2019designing}, amnesic
probing~\cite{elazar2020amnesic}, linear concept erasure~\cite{ravfogel2020null,belrose2023leace} and
direct tests of whether probe accuracy entails task relevance~\cite{ravichander2021probing,ivanova2021probing}
established for language models (\S\ref{sec:decoding}, \S\ref{sec:intervention}). An uncontrolled
intervention that changes behaviour shows causal
sensitivity, not that the direction means what it was named; \S\ref{sec:taxonomy} tags the two apart.

\textbf{Concepts.} A concept here is any property of the input, the report or the acquisition that a human
can name and, in principle, label: a finding or anatomical structure, but also a relational or temporal
property, a report-level attribute, or a nuisance factor such as scanner or site. Features discovered
without a prior vocabulary acquire concept status only through the naming step, which is why provenance is
coded separately from the concept itself. Concepts enter through
five routes (\S\ref{sec:taxonomy}), and the route fixes how validity is checked: label-schema concepts
inherit their annotation's reliability, imperfect for report-extracted radiology
labels~\cite{irvin2019chexpert,smit2020chexbert}; text-mined and model-generated banks scale but need
clinician validation (\S\ref{sec:evaluation}); and discovered features acquire meaning through a
labelling step of studied stability (\S\ref{sec:discovery}).

\subsection{Loci along the computational path}
\label{sec:loci}

Fig.~\ref{fig:overview} places the six loci on the three architectural templates of
\S\ref{sec:models}. A unimodal encoder offers the vision encoder alone; a dual encoder adds a text
encoder and the projected joint space, where zero-shot classification and retrieval take place; a
multimodal LLM adds the connector output and a residual stream at each layer and token position,
written into by attention heads and MLP blocks. The pre-readout representation is the final
hidden state before the unembedding; its logits and tokens are end-points, not loci. A locus is therefore a coordinate $\lambda=(\text{stage},\ell,p,\text{pooling})$, consisting of the stage of the path, the layer, the token or patch position and the pooling applied, and the six labels are its stage-level bins, not disjoint sets of states: the pre-readout representation is the decoder residual stream at its
final layer and target position, and it is kept as a label because the studies that measure it treat it
as a distinct object. Within the vision encoder the coding records the pooling coordinate as a
\emph{sublocus} (pooled or class token, tile, patch, slide aggregate, volume pool), which
\S\ref{sec:findings} shows often decides what is recoverable; a study that repeats one measurement at
several loci is coded once, with the label ``Path'' marking a multi-locus record (\S\ref{sec:taxonomy}).

Two properties recur. First, findings do not transfer across loci: what is decodable from
the encoder need not survive the connector, and what is in the residual stream at one token position
need not influence the logit at another. Second, the language model contributes a prior of its own, so
any measurement after fusion mixes visual evidence with linguistic expectation: the confound is absent
from image-only encoders, attenuated in a dual encoder's joint space and strongest after fusion.

\textbf{Modality-specific structure.} Two-dimensional radiographs and photographs are processed as one
image. Three-dimensional CT and MRI models tokenise volumes or slices and often train their own
encoders~\cite{hamamci2024ctrate,blankemeier2024merlin}, so a ``patch'' is a sub-volume, slice or slab
and pooling to a volume-level representation is a design choice studies rarely state. Whole-slide
pathology models stack an aggregator on frozen tile encoders, so ``the representation'' may
mean a tile or slide embedding or the aggregator's attention
weights~\cite{chen2024generalpurpose,xu2024wholeslide}. Multi-image settings add token positions
carrying temporal information the encoder alone cannot represent~\cite{bannur2024maira2}.
Results on two-dimensional chest radiographs, which dominate the corpus, do not extend automatically.

\subsection{Related surveys}
\label{sec:related}

%% [written after the `surveys' literature bucket is in]
Adjacent surveys form four groups (Table~\ref{tab:surveys}). \emph{Medical XAI surveys} cover attribution
and visualisation, self-explainable concept and prototype designs, clinician-facing explanations and, closest
to this review, mechanistic interpretability in medical report
generation~\cite{tjoa2019survey,vandervelden2021explainable,salahuddin2021transparency,shrestha2023medical,hou2024self,sun2024explainable,imran2026multimodal,muhammad2025cracking}.
\emph{Medical foundation-model and VLM surveys} cover architectures, pre-training, applications, data and
trustworthiness, treating alignment as an objective rather than an analysed
representation~\cite{zhang2023challenges,azad2023foundational,khan2024comprehensive,ye2025multimodal}.
\emph{General interpretability surveys} cover inner structure, mechanisms, probing, concepts, sparse
autoencoders and representation engineering with little medical
content~\cite{rauker2022toward,bereska2024mechanistic,belinkov2021probing,poeta2023concept,shu2025a,wehner2025taxonomy,knab2026whats}.
\emph{Multimodal-LLM interpretability surveys} cover attention and neuron-level methods, hallucination and
misalignment, without medical
models~\cite{lin2025survey,dang2024explainable,shu2025large,bai2024hallucination,zhang2026locate}.

This survey differs in three respects. First, its object is medical vision and multimodal foundation models
and its subject their internal representations, an intersection Table~\ref{tab:surveys} shows to be
covered only partially. Second, it organises the literature by methodology, locus and concept provenance under explicit rules,
coding every core study at the level of study $\times$ family $\times$ evidence type $\times$ locus. Third, it treats evaluation resources and protocol comparability as review
objects, so evidence can be characterised by its breadth and reported safeguards, not only by its existence.

\begin{table*}[!tb]
\centering\small
\caption{Related surveys. Med.: medical scope; FM: foundation models as the object; MM: multimodal
(image--text) models; Rep.: reviews representation-level (internal) analyses; Int.: reviews
interventional evidence (patching, steering, editing); Eval.: consolidates evaluation resources;
Prot.: states a search/inclusion protocol. \cmark\ yes (a dedicated section or explicit inclusion
criterion), \pmark\ partial (mentioned in passing or covered for a subset only), \xmark\ no (absent from
scope, sections and reference list), as read by the authors from each survey's section list and reference
list rather than from a coded rubric, so the borderline \pmark\ cells are a judgement a reader may differ
on; the surveys are cited so each cell can be checked against its source.}
\label{tab:surveys}
\setlength\tabcolsep{5pt}\begin{tabular}{@{}L{5.0cm}cL{2.4cm}ccccccc@{}}
\toprule
Survey & Year & Venue & Med. & FM & MM & Rep. & Int. & Eval. & Prot. \\
\midrule
\multicolumn{10}{@{}l}{\emph{Medical XAI}} \\
Tjoa \& Guan~\cite{tjoa2019survey} & 2021 & TNNLS & \cmark & \xmark & \xmark & \pmark & \xmark & \xmark & \xmark \\
van der Velden \emph{et al.}~\cite{vandervelden2021explainable} & 2022 & MedIA & \cmark & \xmark & \xmark & \xmark & \xmark & \pmark & \cmark \\
Patr\'icio \emph{et al.}~\cite{patricio2022explainable} & 2023 & ACM CSUR & \cmark & \xmark & \pmark & \pmark & \xmark & \cmark & \xmark \\
Hou \emph{et al.}~\cite{hou2024self} & 2024 & arXiv & \cmark & \pmark & \pmark & \pmark & \xmark & \pmark & \cmark \\
Sun \emph{et al.}~\cite{sun2024explainable} & 2025 & ACM THC & \cmark & \xmark & \cmark & \xmark & \xmark & \xmark & \xmark \\
Muhammad \emph{et al.}~\cite{muhammad2025cracking} & 2025 & CSBR & \cmark & \cmark & \cmark & \pmark & \pmark & \xmark & \cmark \\
Imran \& Lee~\cite{imran2026multimodal} & 2026 & IEEE Access & \cmark & \cmark & \cmark & \xmark & \xmark & \pmark & \cmark \\
\midrule
\multicolumn{10}{@{}l}{\emph{Medical foundation models and VLMs}} \\
Zhang \& Metaxas~\cite{zhang2023challenges} & 2023 & MedIA & \cmark & \cmark & \pmark & \xmark & \xmark & \xmark & \xmark \\
Khan \emph{et al.}~\cite{khan2024comprehensive} & 2025 & IEEE RBME & \cmark & \cmark & \cmark & \xmark & \xmark & \xmark & \xmark \\
Zhao \emph{et al.}~\cite{zhao2023clip} & 2025 & MedIA & \cmark & \cmark & \cmark & \xmark & \xmark & \pmark & \xmark \\
Chen \emph{et al.}~\cite{chen2024survey} & 2026 & IJCV & \cmark & \cmark & \cmark & \xmark & \xmark & \pmark & \xmark \\
Xiao \emph{et al.}~\cite{xiao2024comprehensive} & 2024 & Inf.\ Fusion & \cmark & \cmark & \cmark & \xmark & \xmark & \pmark & \xmark \\
Shi \emph{et al.}~\cite{shi2024a} & 2024 & arXiv & \cmark & \cmark & \cmark & \pmark & \xmark & \xmark & \xmark \\
\midrule
\multicolumn{10}{@{}l}{\emph{General interpretability}} \\
R\"auker \emph{et al.}~\cite{rauker2022toward} & 2023 & SaTML & \xmark & \xmark & \xmark & \cmark & \cmark & \pmark & \xmark \\
Bereska \& Gavves~\cite{bereska2024mechanistic} & 2024 & TMLR & \xmark & \pmark & \xmark & \cmark & \cmark & \xmark & \xmark \\
Poeta \emph{et al.}~\cite{poeta2023concept} & 2023 & ACM CSUR & \xmark & \xmark & \xmark & \cmark & \pmark & \cmark & \cmark \\
Shu \emph{et al.}~\cite{shu2025a} & 2025 & EMNLP & \xmark & \cmark & \xmark & \cmark & \pmark & \cmark & \xmark \\
Wan \emph{et al.}~\cite{wan2025survey} & 2026 & TPAMI & \xmark & \pmark & \pmark & \pmark & \pmark & \cmark & \xmark \\
\midrule
\multicolumn{10}{@{}l}{\emph{Multimodal-LLM interpretability}} \\
Lin \emph{et al.}~\cite{lin2025survey} & 2025 & arXiv & \xmark & \cmark & \cmark & \cmark & \cmark & \pmark & \xmark \\
Dang \emph{et al.}~\cite{dang2024explainable} & 2024 & arXiv & \xmark & \cmark & \cmark & \cmark & \pmark & \pmark & \cmark \\
\bottomrule
\end{tabular}
\end{table*}

